\section{Model formulation}\label{sec:formulation}

\subsection{One sequence fixes every start time}\label{sec:timing}

\sumpara A job sequence and the earliest-start recursion together determine the schedule and both measures. \sumparaend
\fin The line processes $n$ jobs, indexed by $j$, on $m$ stages, indexed by $k$ in line order. \fine
\fin Job $j$ is released at time $r_j$, is due at $d_j$, carries weight $w_j$, and needs processing time $p_{jk} > 0$ at stage $k$. \fine
\fin A sequence $\pi$ lists the jobs in the order they enter the line, with $\pi(i)$ the job in position $i$, and the same order holds at every stage. \fine
\fin Given $\pi$, each operation starts at the later of the job's completion at the previous stage and the stage's completion of the previous job: \fine
\begin{equation}\label{eq:recursion}
S_{\pi(i),k} = \max\bigl\{C_{\pi(i),k-1},\; C_{\pi(i-1),k}\bigr\}, \qquad C_{\pi(i),k} = S_{\pi(i),k} + p_{\pi(i),k},
\end{equation}
\fin with the conventions $C_{\pi(i),0} = r_{\pi(i)}$, so the first stage waits for the release, and $C_{\pi(0),k} = -\infty$, so the first job waits for no other. \fine

\sumpara The recursion gives the smallest tardiness any schedule with that sequence can reach, so choosing a schedule reduces to choosing a sequence. \sumparaend
\fin The tardiness of job $j$ is $T_j = \max\{0,\, C_{jm} - d_j\}$, and the two measures of \cref{sec:objective} are \fine
\begin{equation}\label{eq:measures}
Z(\pi) = \sum_{j=1}^{n} w_j T_j, \qquad W(\pi) = \sum_{j=1}^{n} \sum_{k=2}^{m} \bigl(S_{jk} - C_{j,k-1}\bigr).
\end{equation}
\fin By induction over positions and stages, any schedule that keeps sequence $\pi$ and satisfies the constraints of the course project starts each operation no earlier than \eqref{eq:recursion}, so the recursion gives every job its earliest completion at the last stage and $\pi$ its smallest value of $Z$. \fine
\fin The planning problem is therefore to choose the sequence $\pi$ that minimizes $Z(\pi)$, and \eqref{eq:recursion} is also the execution rule that the course project prescribes for the stress test. \fine

\subsection{A mixed-integer program defines the optimum}\label{sec:mip}

\sumpara A positional program states the same problem in a form a solver can certify. \sumparaend
\fin To certify optimal values, we state the problem as a mixed-integer program over sequence positions $q = 1, \ldots, n$. \fine
\fin The binary variable $x_{jq}$ equals 1 when job $j$ takes position $q$, $B_{qk}$ and $E_{qk}$ are the start and end of the job in position $q$ at stage $k$, and $T_j \geq 0$ is the tardiness of job $j$: \fine
\begin{align}
\min \quad & \sum_{j=1}^{n} w_j T_j \label{eq:mip-obj}\\
\text{s.t.} \quad & \textstyle\sum_{q=1}^{n} x_{jq} = 1 \quad \forall j, \qquad \sum_{j=1}^{n} x_{jq} = 1 \quad \forall q, \label{eq:mip-assign}\\
& E_{qk} = B_{qk} + \textstyle\sum_{j=1}^{n} p_{jk}\, x_{jq} \quad \forall q, k, \label{eq:mip-proc}\\
& B_{q1} \geq \textstyle\sum_{j=1}^{n} r_j\, x_{jq} \quad \forall q, \label{eq:mip-release}\\
& B_{qk} \geq E_{q,k-1} \quad \forall q,\ k \geq 2, \label{eq:mip-stage}\\
& B_{qk} \geq E_{q-1,k} \quad \forall q \geq 2,\ k, \label{eq:mip-machine}\\
& T_j \geq E_{qm} - d_j - U_j\,(1 - x_{jq}) \quad \forall j, q, \label{eq:mip-tardy}\\
& x_{jq} \in \{0,1\}, \quad T_j \geq 0. \notag
\end{align}
\fin Constraints \eqref{eq:mip-assign} place every job in exactly one position, \eqref{eq:mip-proc} runs each operation without interruption, \eqref{eq:mip-release} and \eqref{eq:mip-stage} enforce release and stage order, and \eqref{eq:mip-machine} lets a stage start a job only after finishing the previous position, which also keeps one order at every stage. \fine
\fin Constraint \eqref{eq:mip-tardy} charges the tardiness of the job in position $q$ to that job, and it is inactive for every other job when $U_j = \max\{0,\, \max_h r_h + \sum_{h,k} p_{hk} - d_j\}$, with $h$ running over all jobs. \fine

\sumpara The tardiness bound never removes the optimal schedule, so the program's optimum equals the best sequence's tardiness. \sumparaend
\fin Under recursion \eqref{eq:recursion}, every job of every sequence finishes by $\max_h r_h + \sum_{h,k} p_{hk}$, because a completion time is a release plus processing times along one path through the schedule. \fine
\fin The bound $U_j$ therefore never forces a positive tardiness on the earliest-start schedule of any sequence, which by \cref{sec:timing} is the best schedule for that sequence. \fine
\fin The program's optimal value equals $\min_\pi Z(\pi)$, and we use it only as a benchmark, never as one of the two compared methods. \fine

\subsection{The rule sorts once, and iterated greedy searches by insertion}\label{sec:methods}

\sumpara The rule to beat is a single sort by due date. \sumparaend
\fin The earliest due date rule sorts the jobs by $d_j$, breaks ties by job identifier, and times the resulting sequence by \eqref{eq:recursion}, exactly as the course project defines it. \fine

\sumpara Iterated greedy builds a weighted start, then repeatedly destroys and rebuilds the sequence around a local search. \sumparaend
\fin Our researched method is IG\_RS\_LS, the iterated greedy algorithm with local search that Ruiz and St\"utzle proposed for the flow shop with sequence-dependent setup times and the makespan or weighted tardiness objective~\cite{ruiz2008igwt}; our line has no setup times, so we apply its weighted tardiness version with every setup time equal to zero, as shown in \cref{alg:ig}. \fine
\fin It compares sequences by the pair $F(\pi) = (Z(\pi), W(\pi))$ in lexicographic order, and the best position for a job is the position in a partial sequence that gives the smallest $F$, the earliest one on an exact tie. \fine
\fin The start, NEH\_EWDD, orders the jobs by ascending $d_j / w_j$ and inserts them one at a time at their best positions~\cite{ruiz2008igwt}, following the insertion steps of NEH as stated by Ruiz and St\"utzle~\cite{ruiz2007ig}. \fine
\fin Each iteration removes $\delta = 4$ jobs at random, reinserts them in the order removed, and applies a local search that moves single jobs to their best positions until no move improves $F$. \fine
\fin A worse sequence replaces the current one with probability $\exp(-\Delta/\Theta)$, where $\Delta$ is its increase in $Z$ and $\Theta = \lambda \sum_{j,k} p_{jk} / (10\, n\, m)$ with $\lambda = 0.5$, the constant-temperature rule and settings of the source~\cite{ruiz2008igwt}. \fine

\begin{algorithm}[ht]
\caption{Iterated greedy for weighted tardiness, after Ruiz and St\"utzle~\cite{ruiz2008igwt}. $F = (Z, W)$ is compared lexicographically, and every insertion goes to the best position.}
\label{alg:ig}
\small
\begin{algorithmic}[1]
\Require instance; destruction size $\delta = 4$; temperature factor $\lambda = 0.5$; iteration limit $L = 5000$; seed 522
\State $\pi \gets$ NEH\_EWDD: jobs by ascending $d_j / w_j$, each inserted at its best position
\If{$F(\pi^{\mathrm{EDD}}) < F(\pi)$} $\pi \gets \pi^{\mathrm{EDD}}$ \Comment{never start worse than the rule}\EndIf
\State $\pi \gets \textsc{LocalSearch}(\pi)$; \quad $\pi^{*} \gets \pi$; \quad $\Theta \gets \lambda \sum_{j,k} p_{jk} / (10\, n\, m)$
\For{$\ell = 1, \ldots, L$}
  \State remove $\delta$ jobs chosen at random from $\pi$, keeping the removed jobs in order of removal
  \State insert each removed job, in that order, at its best position
  \State $\pi' \gets \textsc{LocalSearch}$ of the rebuilt sequence
  \If{$F(\pi') < F(\pi)$} $\pi \gets \pi'$; \textbf{if} $F(\pi) < F(\pi^{*})$ \textbf{then} $\pi^{*} \gets \pi$
  \ElsIf{$u \leq \exp\bigl(-(Z(\pi') - Z(\pi))/\Theta\bigr)$ for $u$ uniform on $[0,1]$} $\pi \gets \pi'$
  \EndIf
\EndFor
\State \Return $\pi^{*}$
\Statex
\Function{LocalSearch}{$\pi$}
  \Repeat
    \For{each job $j$, in random order}
      \State remove $j$ and insert it at its best position, giving $\pi''$
      \If{$F(\pi'') < F(\pi)$} $\pi \gets \pi''$ \EndIf
    \EndFor
  \Until{a full pass changes nothing}
  \State \Return $\pi$
\EndFunction
\end{algorithmic}
\end{algorithm}

\sumpara Three choices are ours, each made to fit this decision. \sumparaend
\fin We break ties in $Z$ by waiting, because the course project names waiting as the secondary measure and the source compares on one objective. \fine
\fin We start from the earliest due date sequence whenever it beats NEH\_EWDD, so the method can never return a sequence worse than the rule; on all four instances NEH\_EWDD was the better start. \fine
\fin We stop after $L = 5000$ iterations with a fixed seed of 522, where the source stops at a CPU-time limit, because the course requires a rerun on another computer to reproduce the submitted result. \fine
\fin The best-position search times the jobs ahead of each position once and abandons a position as soon as its running pair $(Z, W)$ can no longer beat the best position found, which is safe because $Z$ and $W$ only grow as jobs are appended. \fine
